After having showed the differential performance of \texttt{seine} and naive methods in a simple synthetic example of electoral competition, I moved to an application to real data.

In particular, I applied the semiparametric estimation to the 2020 US presidential election in Texas, using the electoral precincts as the geographical units.


\subsection{Construction of the dataset}

The dataset is constructed by combining electoral results at the electoral precinct level (also referred to as voting districts or VTDs) with demographic data of 2020 US census, with the goal of obtaining, for each precinct, a response variable ($\overline{Y}_g$), a partition of the population in categories ($\overline{X}_g$) and a set of covariates ($Z_g$).

In the construction of the dataset, I exploited three sources:
\begin{itemize}
    \item ALARM dataset by \citet{t._kenny2021}, which combines electoral results at the electoral precinct level (also referred to as voting districts or VTDs) previously collected by \citet{vest2020precinct} with demographic data of 2020 US census collected by \citet{uscensus2020}
    \item American Community Survey (ACS) by the U.S. Census Bureau, which provide socio-economic characteristics of the population and can be called through the \texttt{tidycensus} package in R \citep{tidycensus}
    \item \citet{census2021baf}'s files, which provide the mapping of census blocks to CBGs and VTDs, which is necessary to aggregate the socio-economic characteristics collected at the CBG level in ACS surveys to the VTD level.
\end{itemize}


The construction of the dataset is explained in detail in Appendix B (Section \ref{sec:Appendix_dataset}).

An important note, however, must be made regarding the quality of the data.

In particular, the ALARM dataset by \citet{t._kenny2021} contains a number of precincts with missing electoral results, even if they have relevant voting age population.
At the same time, other precints have a number of votes cast that is higher than the voting age population, which constitutes a clear data error.

In order to overcome these issues, I removed from the dataset all the precincts with less than 10 votes cast (which are either the result of errors or not too relevant for the analysis), and all the precincts with more votes cast than voting age population.
Furthermore, I removed also the precincts where the number of votes cast is less than one fifth of the voting age population, as they looked to be outliers in the distribution of the share of votes cast over voting age population.


\subsection{Choice of variables and validity of the assumptions}

The goal of this application is to estimate the support for the Democratic and Republican candidates in the 2020 US presidential election in Texas among different ethnic groups.

The response variable $\overline{Y}_g$ is the share of votes received by the Democratic candidate in precinct $g$.

The partition of the population $\overline{X}_g$ is the share of voters in precinct $g$ that belong to each of the four ethnic groups considered in the analysis: White, Black, Hispanic and Other.

Finally, the covariate $Z_g$ is the share of 25 years old or older voters in precinct $g$ that have a college degree, which is collected at the CBG level in ACS surveys and then aggregated to the VTD level as explained in Appendix B (Section \ref{sec:Appendix_dataset}).

The assumption of bounded $N$ is satisfied by the distribution of the number of voters across precincts being bounded.

The assumption of positivity is satisfied by the distribution of the share of voters in each ethnic group across precincts with varying levels of college degree holders.

Finally, the assumption of CAR cannot be directly tested, but it is driven by the fact that shares of college degree holders in a precinct are likely to capture many of the unobserved characteristics that affect the voting behavior of different ethnic groups in that precinct, such as location of the county (urban vs rural), income, and other socio-economic characteristics.

\subsection{Results of the estimation}

The results of the estimation are summarized in Table \ref{tab:ei-estimates-texas}, which shows the estimated support for the Democratic (Joe Biden) and Republican (Donald Trump) candidates among each ethnic group, along with the 95\% credible intervals.


% latex table generated in R 4.5.1 by xtable 1.8-4 package
% Wed Sep  9 10:28:30 2026
\begin{table}[ht]
\centering
\begin{tabular}{rllrrrr}
  \hline
 & predictor & outcome & estimate & std.error & conf.low & conf.high \\ 
  \hline
1 & vap\_hisp & pre\_20\_rep\_tru & 0.293 & 0.007 & 0.280 & 0.307 \\ 
  2 & vap\_white & pre\_20\_rep\_tru & 0.854 & 0.010 & 0.834 & 0.875 \\ 
  3 & vap\_black & pre\_20\_rep\_tru & -0.066 & 0.008 & -0.082 & -0.050 \\ 
  4 & other\_ethnicity & pre\_20\_rep\_tru & 0.237 & 0.019 & 0.200 & 0.275 \\ 
  5 & vap\_hisp & pre\_20\_dem\_bid & 0.707 & 0.010 & 0.686 & 0.727 \\ 
  6 & vap\_white & pre\_20\_dem\_bid & 0.146 & 0.005 & 0.136 & 0.156 \\ 
  7 & vap\_black & pre\_20\_dem\_bid & 1.066 & 0.022 & 1.022 & 1.109 \\ 
  8 & other\_ethnicity & pre\_20\_dem\_bid & 0.763 & 0.020 & 0.723 & 0.803 \\ 
   \hline
\end{tabular}
\caption{Semiparametric ecological inference estimates for the 2020 presidential election in Texas} 
\label{tab:ei-estimates-texas}
\end{table}


An evident shortcoming of the results is that the estimates for the support of the candidates among Black voters are not within the [0,1] interval.
That, of course, should cast doubt on the reliability of the estimates in general, but also on the validity of the methodology, at least with regard to this specific application.

In fact, \citet{mccartan2025identificationsemiparametricestimationconditional} explain that, even if the software \texttt{seine} allows to exploit boundedness of response variable in the estimation, it does not guarantee that the estimates will be within the [0,1] interval.


